\subsection[VtfReadStruct]{Read structure file}\label{ssec:VtfReadStruct}

This function reads an input \vsf file line by line, identifying all beads and
molecules present there.

\TODO requires splitting into more functions---I guess\ldots

\CallInit{void}{VtfReadStruct}
\begin{longtable}{L{0.36\textwidth}Mp{0.51\textwidth}}
  \toprule
  variable           & in/out & explanation \\
  \midrule
  \ttb{(char) struct_file[]}         & in  & name of the input \vsf file\\
  \ttb{(bool) detailed}              & in  & should the bead/molecule type
                                             determination be based on more than
                                             just name?\\
  \ttb{(COUNTS) Counts}              & out & basic system information \TODO add
                                             ref\\
  \ttb{(BEADTYPE) BeadType[]}        & out & information about bead types \TODO
                                             add ref\\
  \ttb{(BEAD) Bead[]}                & out & information about individual beads
                                             \TODO add ref\\
  \ttb{(int) Index[]}                & out & array connecting internal and  \vtf
                                             bead indices\\
  \ttb{(MOLECULETYPE) MoleculeType[]}& out & information about molecule types
                                             \TODO add ref\\
  \ttb{(MOLECULE) Molecule[]}       & out & information about individual
                                            molecules \TODO add ref\\
  \bottomrule
\end{longtable}
\vspace{-1em}
\begin{algorithmic}[1]
  \St initialize empty \ttb{Counts}
  \St open \ttb{struct_file}
  \St $file\_line\_count\gets0$\Comment{total number of lines in
    \ttb{struct_file}}
  \Stx $count\_atom\_lines\gets0$\Comment{number of \hltt{a[tom]} lines in
    \ttb{struct_file}}
  \Stx $count\_bond\_lines\gets0$\Comment{number of \hltt{b[ond]} lines in
    \ttb{struct_file}}
  \Stx $default\_atom\_line\gets-1$\Comment{line number of the first
    \hltt{a[tom] default} line}
  \Stx $last\_struct\_line\gets0$\Comment{line number of the last
    \hltt{a[tom]}/\hltt{b[ond]} line}
  \Stx $atom\_names\gets0$\Comment{number of unique bead names in
    \ttb{struct_file}}
  \Stx $res\_names\gets0$\Comment{number of unique molecule names in
    \ttb{struct_file}}
  \Stx $atom\_name[]$ and $resn\_name[]$\Comment{arrays for the unique
    names (memory allocated as needed)}
  \Stx\vspace{-0.6em}\rule{0.96\textwidth}{0.3pt}
  \Stx i) count lines and save bead and molecule names from \ttb{struct_file}
  \Stx\vspace{-0.6em}\rule{0.96\textwidth}{0.3pt}
  \While{$line\gets$ get a line from \ttb{input_vcf}}
    \St $file\_line\_count\gets file\_line\_count+1$
    \If{$line$ is \hltt{a[tom]} line}
      \St $count\_atom\_lines\gets count\_atom\_lines+1$
      \St $last\_struct\_line\gets file\_line\_count$
      \If{$line$ is \hltt{a[tom] default} line}
        \If{not the first \hltt{a[tom] default} line}
          \St warn that this line is not used as the \hltt{default} line
        \Else
          \St $default\_atom\_line\gets file\_line\_count-1$\Comment{$-1$ as
            array indices start from 0}
        \EndIf
      \ElsIf{bead index $>$ \ttb{Counts.BeadsInVsf}}
        \St \ttb{Counts.BeadsInVsf} $\gets$ the bead index$+1$\Comment{$+1$ as
        \vtf bead indices start from 0}
      \EndIf
      \If{bead name is not in $atom\_name[]$}\Comment{tested via a \tt{for}
        loop}
        \St $atom\_names\gets atom\_names+1$
        \St reallocate memory for $atom\_name[]$ and save the name
      \EndIf
      \If{molecule name is not in $res\_name[]$}\Comment{tested via a \tt{for}
        loop}
        \St $res\_names\gets res\_names+1$
        \St reallocate memory for $res\_name[]$ and save the name
      \EndIf
      \If{molecule index $>$ \ttb{Counts.Molecules}}\Comment{count molecules;
        \TODO still assumes continuous numbering starting from 1}
        \St \ttb{Counts.Molecules} $\gets$ the molecule index
      \EndIf
    \ElsIf{$line$ is \hltt{b[ond]} line}
      \St $count\_bond\_lines\gets count\_bond\_lines+1$
      \St $last\_struct\_line\gets file\_line\_count$
    \ElsIf{$line$ is a \hltt{t[imestep]} line}
      \Break
    \ElsIf{$line$ is not recognised}
      \Error unrecognised line
    \EndIf
  \EndWhile

  \Stx\vspace{-0.6em}\rule{0.96\textwidth}{0.3pt}
  \Stx ii) read all structural information from \ttb{struct_file}
  \Stx\vspace{-0.6em}\rule{0.96\textwidth}{0.3pt}

% \St $file\_line\_count \gets 0$\Comment{count lines; line number printed in
%   case of an error}
% \While{$line\gets$ get a line from \ttb{input_vcf}}
%   \St $file\_line\_count\gets file\_line\_count+1$
%   \If{$line$ is \hltt{pbc} line}
%     \St \ttb{Box.Length}$\gets$ 1st to 3rd string from $line$ as box
%       dimensions
%     \St \ttb{Box}$\gets 90^\circ$ as all three angles\Comment{assume
%       orthogonal box}
%     \If{\hltt{pbc} line contains angles}
%       \St \ttb{Box} $\gets$ 4th to 6th string from $line$ as
%         angles\Comment{possibly triclinic box}
%     \EndIf
%     \Break\Comment{box dimensions found, so there's nothing more to do}
%   \ElsIf{$line$ is a coordinate line}
%     \Error missing box dimensions in \ttb{input_vcf}
%   \ElsIf{$line$ is unrecognised}
%     \Error invalid line in \ttb{input_vcf}
%   \EndIf
% \EndWhile
% \St close \ttb{input_vcf}
\end{algorithmic}
\algbottomrule
